\chapter{Podział obowiązków w zespole}

W ramach realizacji projektu prace zostały podzielone pomiędzy członków zespołu w sposób umożliwiający równoległe prowadzenie prac nad różnymi obszarami systemu.

W trakcie realizacji projektu prace były prowadzone w sposób iteracyjny z wykorzystaniem tablicy Kanban. Zadania były dzielone na mniejsze elementy i przypisywane do poszczególnych członków zespołu. Pozwoliło to na bieżące monitorowanie postępu prac, kontrolę priorytetów oraz równoległą realizację poszczególnych modułów systemu.

Poniżej przedstawiono zakres odpowiedzialności poszczególnych osób.

\section{Backend}

Za implementację warstwy backendowej odpowiedzialny był Sebastian Górski.

Zakres obowiązków obejmował:
\begin{itemize}
	\item implementację REST API,
	\item logikę biznesową systemu,
	\item integrację z bazą danych,
	\item implementację mechanizmów autoryzacji i uwierzytelniania.
\end{itemize}

\section{Klient web}

Za implementację aplikacji webowej odpowiedzialny był Jakub Grzymisławski.

Zakres obowiązków obejmował:
\begin{itemize}
	\item implementację interfejsu użytkownika w React,
	\item wykorzystanie TypeScript oraz Vite,
	\item przygotowanie struktury komponentów aplikacji,
	\item implementację nawigacji z użyciem React Router,
	\item integrację z backendem poprzez REST API,
	\item obsługę stanu aplikacji i komunikacji z API.
\end{itemize}

\section{Klient mobile}

Za implementację aplikacji mobilnej odpowiedzialny był Rafał Wilczewski.

Zakres obowiązków obejmował:
\begin{itemize}
	\item implementację aplikacji w React Native,
	\item wykorzystanie Expo SDK,
	\item przygotowanie interfejsu użytkownika w TypeScript,
	\item konfigurację nawigacji z użyciem Expo Router,
	\item integrację z backendem poprzez REST API,
	\item obsługę lokalnego przechowywania danych (Expo Secure Store),
	\item zarządzanie stanem aplikacji.
\end{itemize}

\section{Klient desktop}

Za implementację aplikacji desktopowej odpowiedzialny był Łukasz Szenkiel.

Zakres obowiązków obejmował:
\begin{itemize}
	\item implementację aplikacji w języku Python,
	\item przygotowanie interfejsu użytkownika,
	\item integrację z backendem poprzez REST API,
	\item implementację funkcjonalności panelu administracyjnego,
	\item obsługę podstawowych operacji zarządzania użytkownikami.
\end{itemize}

\section{Bezpieczeństwo systemu}

Za implementację mechanizmów bezpieczeństwa odpowiedzialny był cały zespół.

Zakres obowiązków obejmował:
\begin{itemize}
	\item implementację zabezpieczeń systemów we wszystkich modułach,
	\item skofnigurowanie i opubilkowanie API oraz frontendu z wykrozystaniem protokołu HTTPS i TLS 1.3,
	\item konfigurację 2FA,
	\item poprawnie zaporgramowaną walidację danych wejśćiowych użytkownika,
	\item wdrożenie mechanizmów ochrony przed atakami.
\end{itemize}

\section{Baza danych}

Za projekt i implementację bazy danych odpowiedzialny był:
\begin{itemize}
	\item Sebastian Górski
\end{itemize}

Zakres obowiązków obejmował:
\begin{itemize}
	\item projekt struktury bazy danych,
	\item implementację encji,
	\item konfigurację migracji,
	\item optymalizację relacji między tabelami.
\end{itemize}

\section{Deployment i infrastruktura}

Za przygotowanie środowiska uruchomieniowego odpowiedzialny był:
\begin{itemize}
	\item Sebastian Górski
	\item Jakub Grzymisławski
\end{itemize}

Zakres obowiązków obejmował:
\begin{itemize}
	\item konfigurację Docker oraz Docker Compose,
	\item wdrożenie aplikacji na serwer VPS,
	\item konfigurację HTTPS i certyfikatów TLS,
	\item automatyzację procesu uruchamiania systemu.
\end{itemize}

\section{Testy}

Za implementację testów odpowiedzialny był cały zespół. Każda osoba była odpowiedzialna za przetestowanie własnego modułu.